% !TeX root = main.tex
\section{Introduction}

협력형 다중 에이전트 강화학습(MARL)은 여러 에이전트가 제한된 지역 관측만으로
행동하면서 공동의 목표를 달성하도록 학습한다. QMIX와 같은 가치 분해 방법은
중앙집중식 학습과 분산 실행을 결합해 이 문제를 다루지만~\cite{rashid2018qmix},
보상이 희소하거나 지연되고 에이전트 사이의 행동 의존성이 클수록 유용한 협력 전략을
탐색하기 어렵다. 분산된 에이전트는 작업 할당, 경로 선택, 자원 사용, 위험 회피와
같은 결정을 공동 목표에 맞게 조정해야 한다. 개별 행동과 관측은 짧은 간격으로
바뀌지만 팀 수준의 계획과 서브골은 더 긴 시간 범위에서 의미를 가진다. 수치적
시행착오만으로 이러한 고수준 구조를 발견하려면 많은 상호작용이 필요할 수 있다.

최근 연구는 LLM의 사전 지식과 언어적 추론을 MARL의 고수준 계획에 이용하려 한다.
LEHCA는 환경의 관측 가능한 정보를 구조화된 텍스트로 변환하고, LLM Commander가
전략, 서브골, 의미적 보상 규칙, 행동 제약을 생성하도록 한다. 생성된 지침은 보상
셰이핑과 선택적 행동 제약을 통해 저수준 학습기에 연결된다~\cite{bai2026lehca}.
여기서 의미 접지는 자유 형식의 서브골을 미리 정의된 측정 함수와 행동 규칙에
대응시키는 과정이다. 이 구조는 LLM이 매 스텝 원시 행동을 직접 결정하지 않고,
느린 시간척도에서 탐색의 방향을 제시한다는 장점이 있다.

그러나 이 시간척도의 분리는 새로운 문제를 만든다. LLM 호출에는 추론 지연과 계산
비용이 따르므로 한 번 생성한 지침을 여러 환경 스텝 동안 재사용해야 한다. 그 사이
서브골이 달성되거나 실패할 수 있고, 이용 가능한 자원과 참여 에이전트, 병목과 위험의
위치가 달라질 수 있다. 과업도 탐색, 실행, 복구처럼 서로 다른 조정 국면으로 전환될
수 있다. 처음에는 유용했던 지침도 이런 변화 뒤에는 더 이상 학습 신호를 만들지
못하거나 현재 과업과 맞지 않는 행동을 유도할 수 있다.

기존의 고정 갱신 주기 $F_{\mathrm{update}}$는 단순하고 예측 가능한 기준선이지만,
현재 지침이 무엇을 강조하는지와 실제 환경이 어떻게 변했는지를 사용하지 않는다.
작은 $F_{\mathrm{update}}$는 변화에 빠르게 대응하는 대신 아직 유효한 지침까지
재생성하며 호출 비용과 LLM 출력의 확률적 변동을 늘린다. 큰
$F_{\mathrm{update}}$는 지침을 효율적으로 재사용하지만 필요한 전략 전환을 지연시킬
수 있다. 더욱이 적절한 주기는 환경, 학습 단계, 정책 역량과 지침 내용에 따라 달라질
수 있으므로 하나의 값을 사전에 선택하기 어렵다.

본 연구는 이 문제를 다음과 같이 묻는다.

\begin{quote}
보유 지침에 대응하는 저비용의 실행 신호만으로 지침 갱신이 필요한 시점을 구별하고,
고정 주기보다 효과적인 성능--호출 비용 절충을 얻을 수 있는가?
\end{quote}

RSVP의 핵심 관찰은 언어 지침이 결국 전이에서 계산 가능한 접지 함수와 셰이핑 보상으로
변환된다는 점이다. 따라서 현재 지침이 강조하는 보조 신호가 가까운 미래에 얼마나
발생할지를 예측하면, 별도의 LLM 자기평가 없이 지침과 관련된 변화량을 감시할 수 있다.
RSVP는 접지 함수별 잔여 셰이핑 가치를 예측하고, 발급 상태에 비해 그 값이 누적해서
낮아질 때 Commander를 조기 호출한다. 이 접근의 기여 후보는 다음과 같다.

\begin{itemize}
  \item 고정된 시간 경과가 아니라 현재 지침의 실행 의미에 조건부인 갱신 문제를
        정식화한다.
  \item 기존 접지 함수의 출력을 공유하는 다중 출력 모니터링 모델과 누적 저하
        검출기를 결합해, 추가 LLM 평가 호출 없이 온디맨드 갱신을 구현한다.
  \item 기반 지침의 유용성, 모니터링 모델의 예측력, 검출과 실제 갱신 이득을 분리하는
        단계적 검증 프로토콜을 제시한다.
\end{itemize}

마지막 항목은 특히 중요하다. RSVP는 잘못된 지침을 올바른 지침으로 바꾸는 방법이
아니라, 주어진 지침 생성 및 접지 모듈을 언제 다시 실행할지 결정하는 방법이다.
따라서 기반 지침 채널이 외부 과제에 유용하다는 전제가 먼저 확인되어야 한다.
